\chapter*{Introduction générale}
\addcontentsline{toc}{chapter}{Introduction générale}
\markboth{Introduction générale}{}

Les textes produits en ligne (commentaires, messages, avis) véhiculent
des états affectifs qui dépassent souvent une simple polarité positive ou
négative. La \textbf{détection d'émotions} vise à identifier des catégories
affectives fines (joie, colère, tristesse, peur, etc.), fréquemment
\textbf{plusieurs à la fois} dans un même énoncé. Cette tâche relève de la
classification \emph{multi-étiquettes} : chaque émotion est une décision
binaire, et les cooccurrences ainsi que le \textbf{déséquilibre} des classes
: quelques labels très fréquents, beaucoup d'autres rares, compliquent
l'apprentissage et l'évaluation.

Ce mémoire s'inscrit dans le Mastère de Recherche à l'Institut Supérieur
d'Informatique (ISI), sous l'encadrement de M.~Sahbi Bahroun. Il porte sur
la conception de \textbf{HAKE-MER} (\emph{Hierarchical Attention with
Knowledge Enhancement for Multi-Emotion Recognition}), une architecture
modulaire destinée à mieux exploiter la \textbf{structure} du texte
(phrases, document), les \textbf{interactions entre émotions}, et des
\textbf{ressources sémantiques externes} (lexiques affectifs). L'évaluation s'appuie sur le corpus GoEmotions \cite{demszky2020},
référence pour les émotions fines en anglais, avec un protocole d'ablation à
budget d'entraînement constant et le \textbf{F1-macro} comme métrique
principale. Les campagnes Step~0 à M4 sont exécutées sur DistilBERT-base
(chapitre~\ref{chap:experimentation}) ; RoBERTa ne couvre que le Step~0 flat.

\textbf{Contributions de ce mémoire} :
\begin{itemize}[noitemsep]
\item une architecture en quatre modules cumulables (encodage hiérarchique,
      attention croisée, enrichissement lexical, pondération dynamique) ;
\item un protocole expérimental explicite (hypothèses H1 à H4, deux backbones
      pré-entraînés, comparaison NRC / SenticNet) ;
\item une analyse visant les performances \emph{par classe}, au-delà des
      moyennes globales.
\end{itemize}

Ce rapport enchaîne problématique, état de l'art, conception et
résultats quantitatifs : le chapitre~\ref{chap:experimentation} reporte
les métriques test et tranche H1 à H4.

\paragraph{Organisation du mémoire.}
Le chapitre~\ref{chap:contexte} présente le contexte, la problématique et
les objectifs. Le chapitre~\ref{chap:etat_art} situe HAKE-MER par rapport
à la littérature sur les émotions multiples, l'attention hiérarchique et
l'intégration de connaissances. Le chapitre~\ref{chap:conception} détaille
l'architecture et les hypothèses de recherche. Le
chapitre~\ref{chap:experimentation} fixe le cadre expérimental. Le
chapitre~\ref{chap:perspectives} discute recommandations et prolongements.
Une conclusion générale clôt le document.
